\begin{enumerate}
\item \textbf{Získání kaučuku:} přírodní kaučuk se oddělí z latexu odpařením vody nebo koagulací kyselinou; lze použít také syntetický kaučuk.
\item \textbf{Příprava směsi:} přidají se plniva a vulkanizační činidlo. Hnětení částečně zkracuje řetězce a zvyšuje zpracovatelnost kaučuku.
\item \textbf{Vulkanizace:} při zahřívání se mezi řetězci vytvářejí příčné chemické vazby. Při vulkanizaci sírou vznikají sirné můstky, obvykle při $140$--$150\,^{\circ}\mathrm C$.
\end{enumerate}
Z kaučuku vzniká pružná \textbf{pryž}; množství síry a míra zesítění určují její tvrdost.
